% Standalone note for instance-source and mapping explanation.
\documentclass{article}
\usepackage{amsmath,amssymb}
\usepackage{geometry}
\usepackage{ctex}
\usepackage{booktabs}
\usepackage{hyperref}
\geometry{a4paper, margin=1in}

\begin{document}

\section{算例来源与导入语义}

本文使用的 $10\_6\_2$、$10\_6\_3$、$10\_9\_2$、$10\_9\_3$ 共 16 个小规模算例，
原始文本文件来自 Zhang 等人在 2021 年发表于 \emph{Transportation Science}
的 MVPRP 基准数据。该论文研究的是 \textbf{delivery} 语义下的 multivehicle production routing problem，
并采用 order-up-to (OU) policy：每期先生产，再在期初向客户配送。

\subsection{原始文本字段的含义}

以零售商（论文中的 customer）行为例，文本中的主要字段含义如下：
\begin{center}
\begin{tabular}{ll}
\toprule
字段 & 含义 \\
\midrule
`InitLevel` & 客户在规划期开始时的初始在手库存 \\
`MaxLevel` & 客户库存容量上限 $L_i$ \\
`Demand` & 客户每期需求量（本文数据中为常数需求） \\
`HoldCost` & 客户单位库存持有成本 \\
\bottomrule
\end{tabular}
\end{center}

供应商（论文中的 plant / depot）行中的 `InitLevel`、`MaxLevel`、`ProdCapacity` 分别对应初始库存、
库存容量与单期生产能力。

\subsection{为什么这些算例直接导入后会出现 infeasible}

Zhang 等（2021）的 delivery 模型中，客户首次在第 $t$ 期被配送、且上次访问期为 $v=0$ 时，
论文定义的首次配送量为
\begin{equation}
g^{\text{delivery}}_{i0t}=\sum_{r=1}^{t-1} d_{ir} + (L_i - I_{i0}).
\end{equation}
这意味着：客户 \texttt{InitLevel} 越大，首次访问越可以被推迟。

但本项目中的实现是 \textbf{pickup + production} 语义。
对收集点 $i$，若首次在第 $t$ 期访问，则当前实现采用
\begin{equation}
g^{\text{pickup}}_{i0t}=I_{i0} + \sum_{r=1}^{t} s_{ir},
\end{equation}
其中 $I_{i0}$ 被解释为“期初待回收库存”，$s_{ir}$ 是每期新增可回收量。
如果直接把原始 delivery 数据中的 `InitLevel` 塞进这里的 $I_{i0}$，
就会把“客户期初库存高”误读成“待回收量高”，从而把原本可行的数据压成过早必须访问，
在 $l=6,9$ 的算例上尤其容易导致 infeasible。

\subsection{在不改原始输入文件的前提下的映射方法}

为保持原 benchmark 的“首次访问紧迫度”不变，本文在导入阶段做如下映射：
\begin{align}
s_{it} &:= \texttt{Demand}_i, \qquad t=1,\dots,l,\\
I^{\text{pickup}}_{i0} &:= L_i - \texttt{InitLevel}^{\text{delivery}}_i - \texttt{Demand}_i.
\end{align}
实现中再取
\begin{equation}
I^{\text{pickup}}_{i0} := \max\{0,\; L_i - \texttt{InitLevel}^{\text{delivery}}_i - \texttt{Demand}_i\},
\end{equation}
以避免数值上出现负的初始待回收库存。

此时，pickup 模型中“首次在第 $t$ 期访问不超容量”的条件为
\begin{equation}
I^{\text{pickup}}_{i0} + \sum_{r=1}^{t} s_{ir} \le L_i,
\end{equation}
代入上述映射后可化为
\begin{equation}
\sum_{r=1}^{t-1} \texttt{Demand}_i \le \texttt{InitLevel}^{\text{delivery}}_i.
\end{equation}
也就是说，delivery benchmark 中“客户凭借期初库存最多能撑到哪一期”的信息，
被等价地转写成了 pickup 模型中“收集点最晚何时必须首次访问”的信息。

\subsection{与代码的对应关系}

上述映射已写入
\href{/Users/tangshenghao/IdeaProjects/MVPPRP/src/main/java/instance/Instance.java}{`src/main/java/instance/Instance.java`}：
\begin{itemize}
\item 说明性注释位于文件开头；
\item 零售商初始库存赋值从 `this.Ii0[i] = nd.init;` 改为映射函数；
\item 映射函数定义为 `mapDeliveryInitToPickupInit(...)`。
\end{itemize}

\subsection{关于类型 1/2/3/4 的补充说明}

对本次重算涉及的四组实例 `10\_6\_2`、`10\_6\_3`、`10\_9\_2`、`10\_9\_3`，
`MVPRP1/2/3/4` 在可行性相关参数
（`PeriodNumber`、`VechileNumber`、`VehicleCapacity`、供应商 `InitLevel/MaxLevel/ProdCapacity`、
零售商 `InitLevel/MaxLevel/Demand`）上完全一致。
不同类型主要体现在成本参数或坐标尺度上，因此本次 `Instance.java` 的修正
会一致地影响这 16 个实例的可行性。

\end{document}
